% Preámbulo común para las soluciones de los prácticos de Programación 3
\documentclass[11pt,a4paper]{article}

\usepackage[utf8]{inputenc}
\usepackage[T1]{fontenc}
\usepackage{lmodern}
\usepackage[spanish,es-noshorthands]{babel}
\usepackage[margin=2.3cm]{geometry}
\usepackage{amsmath,amssymb,amsthm}
\usepackage{enumitem}
\usepackage{xcolor}
\usepackage{tikz}
\usetikzlibrary{arrows.meta,positioning,shapes,calc}
\usepackage[noend]{algpseudocode}
\usepackage{algorithm}
\usepackage{booktabs}
\usepackage{tcolorbox}
\usepackage[hidelinks]{hyperref}

\floatname{algorithm}{Algoritmo}
\renewcommand{\algorithmicrequire}{\textbf{Entrada:}}
\renewcommand{\algorithmicensure}{\textbf{Salida:}}
\algrenewcommand\algorithmicwhile{\textbf{mientras}}
\algrenewcommand\algorithmicdo{\textbf{hacer}}
\algrenewcommand\algorithmicfor{\textbf{para}}
\algrenewcommand\algorithmicforall{\textbf{para todo}}
\algrenewcommand\algorithmicif{\textbf{si}}
\algrenewcommand\algorithmicthen{\textbf{entonces}}
\algrenewcommand\algorithmicelse{\textbf{sino}}
\algrenewcommand\algorithmicreturn{\textbf{retornar}}
\algrenewcommand\algorithmicfunction{\textbf{función}}
\algrenewcommand\algorithmicprocedure{\textbf{procedimiento}}

\theoremstyle{plain}
\newtheorem{lema}{Lema}[section]
\newtheorem{teorema}[lema]{Teorema}
\newtheorem{proposicion}[lema]{Proposición}
\theoremstyle{definition}
\newtheorem{definicion}[lema]{Definición}
\theoremstyle{remark}
\newtheorem*{observacion}{Observación}

% Caja para el enunciado de cada ejercicio
\newtcolorbox{enunciado}{colback=gray!6,colframe=gray!50,boxrule=0.4pt,arc=2pt,left=6pt,right=6pt,top=4pt,bottom=4pt,fontupper=\small}
% Nota que vincula el ejercicio con parciales/exámenes
\newcommand{\examref}[1]{\par\noindent{\small\color{blue!50!black}\textbf{Relación con parciales/exámenes:} #1}\par\medskip}

\newcommand{\NP}{\mathsf{NP}}
\newcommand{\Pclass}{\mathsf{P}}
\newcommand{\Oh}{\mathcal{O}}

\setlength{\parskip}{0.4em}
\setlength{\parindent}{0pt}


\tikzset{
  nodo/.style={circle,draw,minimum size=6.5mm,inner sep=1pt,font=\small},
  dir/.style={-{Stealth[length=2.6mm]},thick},
  arbol/.style={very thick},
  noarbol/.style={dashed,gray}
}
\algrenewcommand\algorithmicloop{\textbf{repetir}}
\newcommand{\Adj}{\mathrm{Adj}}
\newcommand{\dist}{\mathrm{dist}}
\newcommand{\Grev}{G^{\mathrm{rev}}}

\title{Prácticos 3 y 4 --- Grafos}
\author{Programación 3 --- Soluciones}
\date{}

\begin{document}
\maketitle

\paragraph{Convenciones.} Salvo que se diga lo contrario, $G=(V,E)$ tiene $n=|V|$ vértices y
$m=|E|$ aristas, es simple (sin lazos ni aristas múltiples) y se representa mediante
\emph{listas de adyacencia} $\Adj[v]$. Usamos los siguientes resultados del libro
(Kleinberg--Tardos, cap.~3), que citamos sin volver a demostrar:
\begin{itemize}[nosep]
  \item \textbf{(3.3)} Las capas $L_0,L_1,\dots$ de BFS desde $s$ cumplen $L_i=\{v:\dist(s,v)=i\}$.
  \item \textbf{(3.4)} En un grafo no dirigido, si $\{x,y\}\in E$, $x\in L_i$, $y\in L_j$, entonces $|i-j|\le 1$.
  \item \textbf{(3.7)} Si $T$ es un árbol DFS y $\{x,y\}\in E\setminus T$, uno de $x,y$ es ancestro del otro en $T$.
  \item \textbf{(3.11), (3.12)} BFS y DFS (y el cálculo de todas las componentes conexas) se implementan en $\Oh(n+m)$ con listas de adyacencia.
\end{itemize}

%======================================================================
\section{Ejercicio 1}
\begin{enunciado}
Una empresa debe instalar $n$ antenas 5G y conectarlas entre sí mediante cables, de modo que haya
conexión entre todas las antenas con la menor cantidad de conexiones.
(a) ¿Cómo modelaría el problema con un grafo?
(b) Alguien propone una solución con $n$ conexiones. ¿Es válida? ¿Es óptima?
(c) Otra persona propone una solución con $n-2$ conexiones. ¿Es válida? ¿Es óptima?
(d) Demuestre que la solución óptima tiene exactamente $n-1$ conexiones.
\end{enunciado}

\subsection*{(a) Modelo}
Cada antena es un vértice: $V=\{1,\dots,n\}$. Cada cable posible entre dos antenas $u,v$ es una
arista no dirigida $\{u,v\}$; como cualquier par de antenas puede cablearse, el conjunto de
conexiones posibles es el de todas las aristas del grafo completo $K_n$. Una \emph{solución} es un
conjunto de aristas $F$ tal que el grafo $H=(V,F)$ es \textbf{conexo} (entre dos antenas cualesquiera
hay un camino de cables, y por lo tanto conexión). Una solución es \emph{óptima} si minimiza $|F|$.
Es decir, buscamos un subgrafo de cubrimiento conexo con la menor cantidad de aristas.

\subsection*{Un lema auxiliar}
\begin{lema}\label{l:componentes}
Todo grafo con $n$ vértices y $k$ aristas tiene al menos $n-k$ componentes conexas.
\end{lema}
\begin{proof}
Inducción en $k$. Si $k=0$ hay exactamente $n$ componentes (vértices aislados).
Sea $H$ con $k\ge1$ aristas, $e=\{u,v\}$ una de ellas y $H'=H-e$, que por hipótesis inductiva tiene
$c'\ge n-(k-1)$ componentes. Al agregar $e$ a $H'$: si $u$ y $v$ están en la misma componente de
$H'$, las componentes no cambian; si están en componentes distintas, esas dos se funden en una y
las demás no cambian. En cualquier caso $H$ tiene $c\ge c'-1\ge n-k$ componentes.
\end{proof}
\begin{observacion}
En consecuencia, un grafo conexo (una sola componente) con $n$ vértices tiene al menos $n-1$ aristas.
\end{observacion}

\subsection*{(b) Solución con $n$ conexiones}
\textbf{Válida:} depende de cuáles sean las $n$ conexiones. Puede ser válida: por ejemplo (para
$n\ge3$) el ciclo $1-2-\dots-n-1$ conecta todas las antenas con $n$ cables. Pero también puede no
serlo: para $n\ge5$, el camino $1-2-\dots-(n-1)$ ($n-2$ aristas) más las aristas $\{1,3\}$ y $\{2,4\}$
suman $n$ aristas distintas y dejan a la antena $n$ aislada. (Para $n\le 4$ cualesquiera $n$ aristas
distintas conectan el grafo, porque un grafo no conexo con $n$ vértices tiene a lo sumo
$\binom{n-1}{2}<n$ aristas; pero la respuesta general es ``no necesariamente''.)

\textbf{Óptima:} nunca. Por la parte (d) existe una solución válida con $n-1<n$ conexiones. Más
concretamente: si $H=(V,F)$ es conexo y $|F|=n$, entonces $H$ tiene un ciclo (un grafo conexo
acíclico es un árbol y tiene $n-1$ aristas); al quitar una arista de un ciclo el grafo sigue siendo
conexo (todo camino que la usaba puede desviarse por el resto del ciclo), y se obtiene una
solución válida con $n-1$ conexiones.

\subsection*{(c) Solución con $n-2$ conexiones}
\textbf{No es válida}, cualesquiera sean las conexiones: por el Lema~\ref{l:componentes} un grafo con
$n$ vértices y $n-2$ aristas tiene al menos $n-(n-2)=2$ componentes conexas, por lo que hay dos
antenas sin conexión entre sí. Al no ser válida, \textbf{tampoco es óptima} (una solución óptima es,
en particular, válida).

\subsection*{(d) La solución óptima tiene exactamente $n-1$ conexiones}
\begin{itemize}
  \item \emph{Cota inferior:} toda solución válida es un grafo conexo con $n$ vértices, por lo que,
  por el Lema~\ref{l:componentes}, tiene al menos $n-1$ aristas.
  \item \emph{Se alcanza:} la estrella $F^*=\{\{1,v\}: 2\le v\le n\}$ tiene $n-1$ aristas y es conexa
  (todo par $u,v$ se conecta por el camino $u-1-v$).
\end{itemize}
Por lo tanto el mínimo de $|F|$ sobre las soluciones válidas es exactamente $n-1$. Además, toda
solución óptima es un árbol: si tuviera un ciclo podríamos quitar una arista de él y seguir conexos
con $n-2$ aristas, contradiciendo la parte (c).

%======================================================================
\section{Ejercicio 2}
\begin{enunciado}
Consideramos el grafo de la figura 2.1.
(a) Escriba 3 algoritmos distintos para recorrer un grafo.
(b) Utilice BFS desde el nodo 8 y muestre el árbol BFS. ¿Es el único árbol BFS posible para este grafo y este nodo inicial?
(c) Utilice DFS desde el nodo 6 y muestre el árbol DFS. ¿Es el único árbol DFS posible?
(d) ¿Existe una recorrida DFS que dé como resultado el árbol de la parte (b)?
(e) ¿Existe una recorrida BFS que dé como resultado el árbol de la parte (c)?
\end{enunciado}

\begin{center}
\begin{tikzpicture}[scale=1.1]
  \node[nodo] (1) at (2,3) {1};  \node[nodo] (7) at (3.5,3) {7};
  \node[nodo] (2) at (1.2,2) {2}; \node[nodo] (3) at (2.8,2) {3};
  \node[nodo] (4) at (0.4,1) {4}; \node[nodo] (5) at (2,1) {5}; \node[nodo] (8) at (3.5,1) {8};
  \node[nodo] (6) at (2,0) {6};
  \draw (1)--(2) (1)--(3) (2)--(3) (2)--(4) (2)--(5) (3)--(5) (3)--(7) (3)--(8) (4)--(5) (5)--(6) (7)--(8);
\end{tikzpicture}

\small Figura 2.1. Listas de adyacencia (en orden creciente):
$1\!:\!2,3$;\; $2\!:\!1,3,4,5$;\; $3\!:\!1,2,5,7,8$;\; $4\!:\!2,5$;\; $5\!:\!2,3,4,6$;\; $6\!:\!5$;\; $7\!:\!3,8$;\; $8\!:\!3,7$.
\end{center}

\subsection*{(a) Tres algoritmos de recorrida}

\begin{algorithm}[H]
\caption{BFS (búsqueda en anchura, por capas)}
\begin{algorithmic}[1]
\Function{BFS}{$G,s$}
  \State $Descubierto[v]\gets\textbf{false}$ para todo $v\in V$; \ $Descubierto[s]\gets\textbf{true}$
  \State $L_0\gets\{s\}$; \ $i\gets 0$; \ $T\gets\emptyset$
  \While{$L_i\neq\emptyset$}
    \State $L_{i+1}\gets\emptyset$
    \ForAll{$u\in L_i$}
      \ForAll{$v\in\Adj[u]$}
        \If{$\neg Descubierto[v]$}
          \State $Descubierto[v]\gets\textbf{true}$; \ $T\gets T\cup\{\{u,v\}\}$; \ agregar $v$ a $L_{i+1}$
        \EndIf
      \EndFor
    \EndFor
    \State $i\gets i+1$
  \EndWhile
  \State \Return $T$
\EndFunction
\end{algorithmic}
\end{algorithm}

\begin{algorithm}[H]
\caption{DFS recursivo (búsqueda en profundidad)}
\begin{algorithmic}[1]
\State Inicialmente $Explorado[v]\gets\textbf{false}$ para todo $v$; se invoca \Call{DFS}{$s$}.
\Function{DFS}{$u$}
  \State $Explorado[u]\gets\textbf{true}$
  \ForAll{$v\in\Adj[u]$}
    \If{$\neg Explorado[v]$}
      \State $T\gets T\cup\{\{u,v\}\}$; \ \Call{DFS}{$v$}
    \EndIf
  \EndFor
\EndFunction
\end{algorithmic}
\end{algorithm}

\begin{algorithm}[H]
\caption{Exploración genérica de la componente de $s$ (libro, sección 3.2)}
\begin{algorithmic}[1]
\Function{Componente}{$G,s$}
  \State $R\gets\{s\}$; \ $T\gets\emptyset$
  \While{existe una arista $\{u,v\}$ con $u\in R$, $v\notin R$}
    \State $R\gets R\cup\{v\}$; \ $T\gets T\cup\{\{u,v\}\}$
  \EndWhile
  \State \Return $R$, $T$
\EndFunction
\end{algorithmic}
\end{algorithm}
Los tres recorren exactamente la componente conexa de $s$; BFS y DFS son dos maneras concretas de
elegir la arista $\{u,v\}$ del algoritmo genérico (también vale como tercer algoritmo el DFS
iterativo con pila del Ejercicio~3).

\subsection*{(b) BFS desde 8}
Usando las listas en orden creciente:
\begin{itemize}[nosep]
  \item $L_0=\{8\}$.
  \item Se procesa 8: descubre 3 y 7. \ $L_1=\{3,7\}$.
  \item Se procesa 3: descubre 1, 2, 5 (7 y 8 ya descubiertos). Se procesa 7: nada nuevo. \ $L_2=\{1,2,5\}$.
  \item Se procesa 1: nada. Se procesa 2: descubre 4. Se procesa 5: descubre 6 (4 ya descubierto). \ $L_3=\{4,6\}$.
  \item Se procesan 4 y 6: nada nuevo; $L_4=\emptyset$.
\end{itemize}
Árbol BFS $T_b=\{\{8,3\},\{8,7\},\{3,1\},\{3,2\},\{3,5\},\{2,4\},\{5,6\}\}$:
\begin{center}
\begin{tikzpicture}[level distance=11mm,
  level 1/.style={sibling distance=30mm},
  level 2/.style={sibling distance=14mm},
  level 3/.style={sibling distance=10mm},
  every node/.style={nodo}]
\node {8}
  child {node {3}
    child {node {1}}
    child {node {2} child {node {4}}}
    child {node {5} child {node {6}}}}
  child {node {7}};
\end{tikzpicture}
\end{center}
\textbf{No es el único.} Las capas sí son únicas (por (3.3), $L_i$ es el conjunto de vértices a
distancia $i$ de 8), pero el padre de un vértice $v\in L_{i+1}$ es el primer vértice de $L_i$
procesado que es adyacente a $v$, y eso depende del orden de las listas. El vértice 4 tiene dos
vecinos en $L_2$, el 2 y el 5: si en $L_2$ se procesa 5 antes que 2 (por ejemplo, si la lista de 3 es
$5,2,1,7,8$), el padre de 4 es 5 y se obtiene el árbol
$T_b'=(T_b\setminus\{\{2,4\}\})\cup\{\{5,4\}\}\neq T_b$. Los demás vértices tienen un único vecino en
la capa anterior (1, 2 y 5 solo al 3; 3 y 7 solo al 8; 6 solo al 5), así que existen exactamente
dos árboles BFS desde 8: $T_b$ y $T_b'$.

\subsection*{(c) DFS desde 6}
Con las listas en orden creciente: \textsc{DFS}(6) $\to$ \textsc{DFS}(5) (primer vecino de 5 no
explorado: 2) $\to$ \textsc{DFS}(2) $\to$ \textsc{DFS}(1) $\to$ \textsc{DFS}(3) (1 y 2 explorados; el
siguiente es 5, explorado; luego 7) $\to$ \textsc{DFS}(7) $\to$ \textsc{DFS}(8). Desde 8 no hay
nada nuevo; se retrocede hasta 2, donde queda 4 sin explorar: \textsc{DFS}(4). Fin.

Árbol DFS $T_c=\{\{6,5\},\{5,2\},\{2,1\},\{1,3\},\{3,7\},\{7,8\},\{2,4\}\}$:
\begin{center}
\begin{tikzpicture}[level distance=9mm,
  level 3/.style={sibling distance=16mm},
  every node/.style={nodo}]
\node {6}
  child {node {5}
    child {node {2}
      child {node {1} child {node {3} child {node {7} child {node {8}}}}}
      child {node {4}}}};
\end{tikzpicture}
\end{center}
\textbf{No es el único.} El árbol depende del orden en que se recorren las listas de adyacencia.
Por ejemplo, si la lista de 5 empieza por 4, se obtiene \textsc{DFS}(6) $\to$ 5 $\to$ 4 $\to$ 2
$\to$ 1 $\to$ 3 $\to$ 7 $\to$ 8, es decir el camino $6-5-4-2-1-3-7-8$, que es un árbol distinto
de $T_c$ (contiene la arista $\{5,4\}\notin T_c$).

\subsection*{(d) ¿Hay una recorrida DFS que produzca $T_b$?}
\examref{Examen diciembre 2017, Pregunta 1: si DFS y BFS desde el mismo vértice $u$ producen el mismo árbol $T$, entonces $G=T$.}
\textbf{Desde el mismo nodo inicial 8, no.} Las aristas de $G$ que no están en $T_b$ son
$\{1,2\}$, $\{2,5\}$, $\{3,7\}$, $\{4,5\}$. En $T_b$ con raíz 8, los vértices 3 y 7 son hermanos (ninguno es
ancestro del otro), y sin embargo $\{3,7\}\in E\setminus T_b$. Por (3.7) $T_b$ no puede ser un árbol
DFS con raíz 8. (También se deduce del resultado del examen citado: como $G\neq T_b$, un DFS desde 8
no puede producir el mismo árbol que el BFS desde 8.)

\textbf{Si se permite empezar en otro nodo, sí.} Busquemos una raíz $r$ tal que, en $T_b$ enraizado en
$r$, cada arista no arbórea una un ancestro con un descendiente. Para una arista no arbórea
$\{x,y\}$ esto ocurre si y solo si $r$ está ``del lado de $x$'' o ``del lado de $y$'' del camino de
$T_b$ entre $x$ e $y$:
\begin{itemize}[nosep]
  \item $\{1,2\}$ (camino $1-3-2$): $r\in\{1\}\cup\{2,4\}$.
  \item $\{2,5\}$ (camino $2-3-5$): $r\in\{2,4\}\cup\{5,6\}$.
  \item $\{4,5\}$ (camino $4-2-3-5$): $r\in\{4\}\cup\{5,6\}$.
  \item $\{3,7\}$ (camino $3-8-7$): $r\in\{1,2,3,4,5,6\}\cup\{7\}$.
\end{itemize}
La única raíz posible es $r=4$, y efectivamente funciona: \textsc{DFS}(4) con listas
$4\!:\!2,5$;\ $2\!:\!3,1,4,5$;\ $3\!:\!1,5,8,2,7$;\ $5\!:\!6,2,3,4$;\ $8\!:\!7,3$ recorre
$4\to2\to3\to1$ (vuelve), $3\to5\to6$ (vuelve), $3\to8\to7$, y produce exactamente las aristas
$\{4,2\},\{2,3\},\{3,1\},\{3,5\},\{5,6\},\{3,8\},\{8,7\}$, que son las de $T_b$.

\subsection*{(e) ¿Hay una recorrida BFS que produzca $T_c$?}
\textbf{No, desde ningún nodo inicial.} Si $T$ es un árbol BFS con raíz $r$, por (3.3) la
profundidad de cada vértice en $T$ es su distancia a $r$ en $G$; y como dos vértices adyacentes
están a distancias que difieren en a lo sumo 1, \emph{toda arista de $G$ une vértices cuyas
profundidades en $T$ difieren en a lo sumo 1}.

Desde 6 (el mismo inicio de la parte (c)) esto falla de inmediato: en $T_c$ el vértice 3 tiene
profundidad 4 (camino $6-5-2-1-3$) pero $\dist(6,3)=2$ ($6-5-3$).

Para una raíz $r$ cualquiera: la arista $\{3,8\}\notin T_c$ y el camino de $T_c$ entre 3 y 8 es
$3-7-8$. Si $r\neq 7$, la raíz está del lado de 3 o del lado de 8 de ese camino, y entonces uno de
3, 8 está a profundidad 2 más que el otro, absurdo. Luego $r=7$; pero con raíz 7 el vértice 2 tiene
profundidad 3 en $T_c$ (camino $7-3-1-2$) mientras que $\dist(7,2)=2$ ($7-3-2$), absurdo. Por lo
tanto ninguna recorrida BFS produce $T_c$.

%======================================================================
\section{Ejercicio 3}
\begin{enunciado}
(a) Implemente el algoritmo DFS de forma iterativa utilizando una pila.
(b) ¿El árbol de recorrida DFS depende de la implementación? Justifique.
(c) ¿Qué modificaciones se le deben realizar al algoritmo DFS para que retorne el árbol de recorrida?
\end{enunciado}

\subsection*{(a) DFS iterativo con pila}
\begin{algorithm}[H]
\caption{DFS iterativo (versión del libro)}
\begin{algorithmic}[1]
\Function{DFS-Iter}{$G,s$}
  \State $Explorado[v]\gets\textbf{false}$ para todo $v\in V$
  \State $S\gets$ pila vacía; \ \Call{Push}{$S,s$}
  \While{$S$ no es vacía}
    \State $u\gets$ \Call{Pop}{$S$}
    \If{$\neg Explorado[u]$}
      \State $Explorado[u]\gets\textbf{true}$
      \ForAll{$v\in\Adj[u]$}
        \State \Call{Push}{$S,v$}
      \EndFor
    \EndIf
  \EndWhile
\EndFunction
\end{algorithmic}
\end{algorithm}
Cada vértice se marca a lo sumo una vez y, al marcarse $u$, se apilan sus $\deg(u)$ vecinos; en
total se hacen $\sum_u\deg(u)=2m$ apilados (más el inicial), por lo que el algoritmo es $\Oh(n+m)$.

\subsection*{(b) ¿El árbol depende de la implementación?}
\textbf{Sí.} Tanto la versión recursiva como la iterativa producen \emph{algún} árbol DFS (ambos
cumplen la propiedad (3.7)), pero no necesariamente el mismo: el árbol depende del orden en que se
consideran los vecinos, y la pila invierte ese orden (el último vecino apilado es el primero en
explorarse). Ejemplo: en el grafo de la figura 2.1 desde 6, con las listas en orden creciente, la
versión recursiva produce $T_c$ (Ejercicio~2(c)). La versión iterativa hace: explora 6, apila 5;
explora 5 (padre 6), apila $2,3,4,6$; desapila 6 (ya explorado) y explora 4 (padre 5), apila $2,5$;
desapila 5 (explorado) y explora 2 (padre 4), apila $1,3,4,5$; descarta 5, 4 y explora 3 (padre 2),
apila $1,2,5,7,8$; explora 8 (padre 3), apila $3,7$; explora 7 (padre 8); descarta y finalmente
explora 1 (padre 3). Árbol: $\{6,5\},\{5,4\},\{4,2\},\{2,3\},\{3,8\},\{8,7\},\{3,1\}$, distinto de
$T_c$ (por ejemplo $\{5,4\}\notin T_c$).

\subsection*{(c) Modificaciones para retornar el árbol}
Hay que recordar, para cada vértice $v\neq s$, desde qué vértice fue descubierto ($padre[v]$); el
árbol es $T=\{\{padre[v],v\}: v\neq s \text{ explorado}\}$.
\begin{itemize}
  \item \emph{Versión recursiva:} al invocar \textsc{DFS}($v$) desde \textsc{DFS}($u$) se agrega
  $\{u,v\}$ a $T$ (como en el Algoritmo~2 del Ejercicio~2).
  \item \emph{Versión iterativa:} se apilan pares (vértice, quién lo apiló) y la arista se agrega
  recién cuando el vértice se \emph{explora} (no cuando se apila, porque un vértice puede apilarse
  muchas veces):
\end{itemize}
\begin{algorithm}[H]
\caption{DFS iterativo que retorna el árbol}
\begin{algorithmic}[1]
\Function{DFS-Arbol}{$G,s$}
  \State $Explorado[v]\gets\textbf{false}$ para todo $v$; \ $T\gets\emptyset$
  \State $S\gets$ pila vacía; \ \Call{Push}{$S,(s,\textbf{nil})$}
  \While{$S$ no es vacía}
    \State $(u,p)\gets$ \Call{Pop}{$S$}
    \If{$\neg Explorado[u]$}
      \State $Explorado[u]\gets\textbf{true}$; \ $padre[u]\gets p$
      \If{$p\neq\textbf{nil}$} $T\gets T\cup\{\{p,u\}\}$ \EndIf
      \ForAll{$v\in\Adj[u]$}
        \If{$\neg Explorado[v]$} \Call{Push}{$S,(v,u)$} \EndIf
      \EndFor
    \EndIf
  \EndWhile
  \State \Return $T$
\EndFunction
\end{algorithmic}
\end{algorithm}
El costo sigue siendo $\Oh(n+m)$: solo se agregan operaciones $\Oh(1)$ por apilado o exploración.

%======================================================================
\section{Ejercicio 4}
\begin{enunciado}
(K\&T 3.2) Dé un algoritmo para detectar si un grafo (no dirigido) $G$ contiene un ciclo; si la
respuesta es afirmativa, debe mostrar un ciclo como evidencia. Debe admitir una implementación
$\Oh(m+n)$.
(a) Presente un algoritmo basado en DFS (reescriba los algoritmos del libro que utilice).
(b) Demuestre su corrección: responde afirmativamente sii $G$ tiene un ciclo, y en ese caso la
evidencia es un ciclo de $G$.
(c) Demuestre que admite una implementación $\Oh(m+n)$.
(d) Explique si es posible resolver el problema sin basarse en DFS.
\end{enunciado}
\examref{Examen julio 2015, Problema 2 (parte 1): obtener un ciclo de un grafo conexo con DFS y un arreglo de padres, reconstruyendo el ciclo subiendo por los padres. Aquí generalizamos a grafos no necesariamente conexos.}

\subsection*{(a) Algoritmo}
Idea: hacemos un DFS completo (sobre todas las componentes) guardando $padre$ y profundidad
$prof$ de cada vértice. Una arista $\{x,y\}$ es \emph{de árbol} si $padre[x]=y$ o $padre[y]=x$. Si
hay alguna arista que no es de árbol, por (3.7) uno de sus extremos es ancestro del otro y el camino
del árbol entre ellos más esa arista es un ciclo.

\begin{algorithm}[H]
\caption{Detección de ciclo con evidencia}
\begin{algorithmic}[1]
\Function{BuscarCiclo}{$G$}
  \State $Explorado[v]\gets\textbf{false}$, \ $padre[v]\gets\textbf{nil}$ para todo $v\in V$
  \ForAll{$s\in V$}
    \If{$\neg Explorado[s]$} \State $prof[s]\gets 0$; \ \Call{DFS}{$s$} \EndIf
  \EndFor
  \ForAll{$x\in V$} \Comment{recorre cada arista $\{x,y\}$ (dos veces)}
    \ForAll{$y\in\Adj[x]$}
      \If{$padre[x]\neq y$ \textbf{and} $padre[y]\neq x$} \Comment{arista no arbórea}
        \If{$prof[x]>prof[y]$} intercambiar $x$ e $y$ \EndIf
        \State $C\gets\langle y\rangle$; \ $w\gets y$
        \While{$w\neq x$}
          \State $w\gets padre[w]$; \ agregar $w$ al final de $C$
        \EndWhile
        \State \Return ``Sí'', $C$ \Comment{el ciclo es $C$ cerrado con la arista $\{x,y\}$}
      \EndIf
    \EndFor
  \EndFor
  \State \Return ``No''
\EndFunction
\Statex
\Function{DFS}{$u$}
  \State $Explorado[u]\gets\textbf{true}$
  \ForAll{$v\in\Adj[u]$}
    \If{$\neg Explorado[v]$}
      \State $padre[v]\gets u$; \ $prof[v]\gets prof[u]+1$; \ \Call{DFS}{$v$}
    \EndIf
  \EndFor
\EndFunction
\end{algorithmic}
\end{algorithm}

\subsection*{(b) Corrección}
Sea $F=\{\{v,padre[v]\}: padre[v]\neq\textbf{nil}\}$ el conjunto de aristas de árbol. Por
construcción $F$ es la unión de los árboles DFS de cada componente (el bosque DFS), y
$prof[v]=prof[padre[v]]+1$ para toda arista de $F$. Como $G$ es simple, una arista $\{x,y\}$ está en
$F$ si y solo si $padre[x]=y$ o $padre[y]=x$, que es exactamente lo que testea la línea 8.

\paragraph{Si responde ``Sí'', $C$ es un ciclo.} Se encontró $\{x,y\}\in E\setminus F$. Ambos
extremos están en la misma componente, luego en el mismo árbol DFS $T$, y por (3.7) uno es ancestro
del otro. Tras el intercambio $prof[x]\le prof[y]$; como $x\ne y$ (no hay lazos) y un ancestro
propio tiene profundidad estrictamente menor, $x$ es ancestro propio de $y$ y $prof[x]<prof[y]$.
Entonces subiendo por $padre$ desde $y$ se llega a $x$ en $k=prof[y]-prof[x]$ pasos (la profundidad
baja de a 1 y el único ancestro de $y$ a profundidad $prof[x]$ es $x$); el ciclo \emph{while}
termina y $C=\langle y=w_0,w_1,\dots,w_k=x\rangle$, donde cada $\{w_i,w_{i+1}\}\in F\subseteq E$. Los
$w_i$ son distintos (tienen profundidades distintas). Además $k\ge 2$: si $k=1$ entonces $x=padre[y]$
y $\{x,y\}\in F$, contradicción. Así, $C$ tiene al menos 3 vértices distintos, aristas consecutivas
en $E$ y la arista de cierre $\{w_k,w_0\}=\{x,y\}\in E$: es un ciclo de $G$.

\paragraph{Si $G$ tiene un ciclo, responde ``Sí''.} El algoritmo responde ``No'' solo si todas las
aristas están en $F$, o sea $E=F$. Mostremos que $F$ no tiene ciclos, de donde si $G$ tiene un ciclo
entonces $E\neq F$, se encuentra una arista no arbórea y la respuesta es ``Sí''. Supongamos que
$v_1,\dots,v_k,v_1$ es un ciclo con aristas en $F$, y sea $v_j$ un vértice del ciclo de profundidad
máxima. Sus dos vecinos en el ciclo, $v_{j-1}$ y $v_{j+1}$ (índices módulo $k$, distintos por ser
$k\ge3$), están unidos a $v_j$ por aristas de $F$; cada arista de $F$ une un vértice con su padre, que
tiene profundidad una unidad menor. Como ninguno de los dos vecinos tiene profundidad mayor que
$v_j$, ninguno es hijo de $v_j$; luego ambos son $padre[v_j]$, es decir $v_{j-1}=v_{j+1}$, absurdo.

\subsection*{(c) Tiempo de ejecución}
Con listas de adyacencia:
\begin{itemize}[nosep]
  \item Inicialización: $\Oh(n)$.
  \item Los DFS de las líneas 3--5: es el cálculo de todas las componentes conexas; cada \textsc{DFS}($u$)
  se invoca una única vez por vértice (se invoca solo si $u$ no está explorado y lo marca enseguida) y
  recorre $\Adj[u]$ con trabajo $\Oh(1)$ por elemento. Total $\Oh(\sum_u(1+\deg u))=\Oh(n+2m)=\Oh(n+m)$.
  \item El doble \emph{para} de las líneas 6--7 recorre todas las listas: $\Oh(n+m)$, con test $\Oh(1)$.
  \item La reconstrucción del ciclo se hace \emph{una sola vez} y sube a lo sumo $n-1$ niveles: $\Oh(n)$.
\end{itemize}
Por la regla de la suma, el total es $\Oh(n+m)$.

\subsection*{(d) ¿Sin DFS?}
Sí. Por ejemplo con BFS: se calcula el bosque BFS $F$ (BFS desde cada vértice no descubierto, guardando
$padre$ y nivel), y se busca una arista $\{x,y\}\notin F$. El mismo argumento de (b) muestra que $F$
es acíclico, así que $G$ tiene un ciclo si y solo si existe tal arista. En ese caso $x$ e $y$ están en
el mismo árbol BFS, y por (3.4) sus niveles difieren en a lo sumo 1; subiendo por $padre$ desde el más
profundo hasta igualar niveles y luego desde ambos simultáneamente hasta que coincidan, se halla su
ancestro común más bajo $z$. Los caminos $x\leadsto z$ e $y\leadsto z$ en el árbol solo comparten $z$, y
junto con $\{x,y\}$ forman un ciclo (de al menos 3 vértices, pues $\{x,y\}$ no es arista de árbol).
Todo es $\Oh(n+m)$. (Para decidir solo la existencia, sin evidencia, alcanza con contar: $G$ es acíclico
si y solo si $m=n-c$, donde $c$ es la cantidad de componentes conexas.)

%======================================================================
\section{Ejercicio 5}
\begin{enunciado}
(K\&T 3.9) $G=(V,E)$ conexo con $n$ vértices y $m$ aristas, y $s,t$ con $\dist(s,t)>n/2$.
(a) Demuestre que existe un nodo $v$ tal que al eliminarlo $s$ y $t$ quedan en componentes conexas
distintas. \emph{Sugerencia:} sea $L_i$ el conjunto de nodos a distancia $i$ de $s$; muestre que
existe $i$, $1\le i\le n/2$, con $|L_i|=1$.
(b) Presente un algoritmo que encuentra dicho nodo en $\Oh(m+n)$.
(c) Muestre que su algoritmo admite una implementación $\Oh(m+n)$.
\end{enunciado}

\subsection*{(a) Existencia del nodo}
Sea $d=\dist(s,t)\ge\lfloor n/2\rfloor+1$ y $L_i=\{v:\dist(s,v)=i\}$ (las capas de BFS desde $s$, por
(3.3)). Sea $k=\lfloor n/2\rfloor$.

\emph{Cada capa $L_i$ con $0\le i\le d$ es no vacía:} si $s=u_0,u_1,\dots,u_d=t$ es un camino mínimo,
entonces $\dist(s,u_i)=i$ (el prefijo hasta $u_i$ es un camino mínimo, pues si hubiera uno más corto
habría uno más corto hasta $t$), así que $u_i\in L_i$.

\emph{Alguna capa intermedia tiene un solo nodo:} las capas $L_0=\{s\}$, $L_1,\dots,L_k$ y $L_d\ni t$
son disjuntas, y $1\le i\le k<d$ para las intermedias. Si fuera $|L_i|\ge2$ para todo $1\le i\le k$,
tendríamos
\[ n\ \ge\ |L_0|+\sum_{i=1}^{k}|L_i|+|L_d|\ \ge\ 1+2k+1\ \ge\ 1+(n-1)+1\ =\ n+1, \]
usando $2\lfloor n/2\rfloor\ge n-1$. Absurdo. Luego existe $i$ con $1\le i\le k\le n/2$ y $L_i=\{v\}$.
Notar que $v\ne s$ (pues $i\ge1$) y $v\neq t$ (pues $i<d$).

\emph{Eliminar $v$ separa a $s$ de $t$:} sea $P=x_0,x_1,\dots,x_r$ un camino de $s$ a $t$ en $G$ y sea
$\ell(x)$ la capa de $x$. Por (3.4), $|\ell(x_{j+1})-\ell(x_j)|\le1$, con $\ell(x_0)=0<i<d=\ell(x_r)$.
Como la capa avanza de a lo sumo una unidad por paso, algún $x_j$ cumple $\ell(x_j)=i$ (sea $j$ el
primer índice con $\ell(x_j)\ge i$; entonces $\ell(x_{j-1})\le i-1$ y $\ell(x_j)\le\ell(x_{j-1})+1\le
i$). Es decir $x_j\in L_i=\{v\}$: \emph{todo} camino de $s$ a $t$ pasa por $v$, así que en
$G-v$ no hay camino de $s$ a $t$, y quedan en componentes conexas distintas.

\subsection*{(b) Algoritmo}
\begin{algorithm}[H]
\caption{Vértice que separa $s$ de $t$}
\begin{algorithmic}[1]
\Function{Separador}{$G,s,t$}
  \State $Descubierto[v]\gets\textbf{false}$ para todo $v$; \ $Descubierto[s]\gets\textbf{true}$
  \State $L_0\gets\{s\}$; \ $i\gets0$; \ $cant[0]\gets1$
  \While{$L_i\neq\emptyset$} \Comment{BFS del libro, guardando el tamaño de cada capa}
    \State $L_{i+1}\gets\emptyset$; \ $cant[i+1]\gets 0$
    \ForAll{$u\in L_i$}
      \ForAll{$w\in\Adj[u]$}
        \If{$\neg Descubierto[w]$}
          \State $Descubierto[w]\gets\textbf{true}$; \ $capa[w]\gets i+1$
          \State agregar $w$ a $L_{i+1}$; \ $cant[i+1]\gets cant[i+1]+1$
        \EndIf
      \EndFor
    \EndFor
    \State $i\gets i+1$
  \EndWhile
  \State $d\gets capa[t]$
  \For{$i\gets1$ \textbf{hasta} $d-1$}
    \If{$cant[i]=1$} \State \Return el único elemento de $L_i$ \EndIf
  \EndFor
\EndFunction
\end{algorithmic}
\end{algorithm}
\emph{Corrección:} por (3.3) las listas $L_i$ calculadas son las capas de distancia y $capa[t]=d$. Por
la parte (a) existe $i\in[1,d-1]$ con $|L_i|=1$, así que el \emph{para} de la línea 13 siempre retorna,
y lo que retorna es un $v$ con $L_i=\{v\}$, $1\le i\le d-1$; la parte (a) muestra que para
\emph{cualquier} tal $i$ (no solo para $i\le n/2$) el vértice $v$ separa a $s$ de $t$, pues el argumento
del tercer párrafo solo usó $0<i<d$.

\subsection*{(c) Tiempo de ejecución}
Las líneas 2--12 son BFS con una operación $\Oh(1)$ extra por vértice descubierto: $\Oh(n+m)$ por
(3.11) (cada vértice entra a una sola capa, y cada lista $\Adj[u]$ se recorre una sola vez, cuando $u$
se procesa). Representando cada $L_i$ como lista enlazada y $cant$ como arreglo, el bucle final
hace a lo sumo $d-1<n$ iteraciones de costo $\Oh(1)$. Total: $\Oh(n+m)$.

%======================================================================
\section{Ejercicio 6}
\begin{enunciado}
(K\&T 3.10) Dado un grafo $G$ y vértices $u,v$:
(a) Presente un algoritmo que calcula la cantidad de caminos de largo mínimo entre $u$ y $v$, en
$\Oh(m+n)$. \emph{Sugerencia:} calcule para cada $w$ la cantidad $N_w$ de caminos de largo mínimo de
$u$ a $w$.
(b) Demuestre la corrección. \emph{Sugerencia:} escriba una recurrencia para $N_w$ en función de los
$N_t$ de los vértices $t$ que ocurren inmediatamente antes que $w$ en un camino mínimo.
(c) Demuestre que admite una implementación $\Oh(m+n)$.
\end{enunciado}
\examref{Examen julio 2014, Problema 2(b) (mismo problema, resuelto con BFS acumulando cantidades); Examen julio 2019, Ej.~2 (cantidad de historias de largo mínimo: la misma recurrencia $N(w)=\sum N(t)$ sobre los $t$ del nivel BFS anterior).}

\subsection*{(a) Algoritmo}
Hacemos BFS desde $u$ por capas. Cuando se procesa $t\in L_i$ y se mira la arista $(t,w)$: si $w$ no
estaba descubierto, entra a $L_{i+1}$ con $N[w]\gets N[t]$; si ya estaba descubierto y está en
$L_{i+1}$, se suma $N[w]\gets N[w]+N[t]$; en otro caso no se hace nada.

\begin{algorithm}[H]
\caption{Cantidad de caminos mínimos de $u$ a $v$}
\begin{algorithmic}[1]
\Function{CantCaminosMin}{$G,u,v$}
  \State $Descubierto[w]\gets\textbf{false}$, \ $N[w]\gets0$ para todo $w\in V$
  \State $Descubierto[u]\gets\textbf{true}$; \ $capa[u]\gets0$; \ $N[u]\gets1$
  \State $L_0\gets\{u\}$; \ $i\gets0$
  \While{$L_i\neq\emptyset$}
    \State $L_{i+1}\gets\emptyset$
    \ForAll{$t\in L_i$}
      \ForAll{$w\in\Adj[t]$}
        \If{$\neg Descubierto[w]$}
          \State $Descubierto[w]\gets\textbf{true}$; \ $capa[w]\gets i+1$; \ agregar $w$ a $L_{i+1}$
        \EndIf
        \If{$capa[w]=i+1$}
          \State $N[w]\gets N[w]+N[t]$
        \EndIf
      \EndFor
    \EndFor
    \State $i\gets i+1$
  \EndWhile
  \State \Return $N[v]$ \Comment{$0$ si $v$ no es alcanzable desde $u$}
\EndFunction
\end{algorithmic}
\end{algorithm}
(La condición de la línea 11 se evalúa solo si $w$ está descubierto, por lo que $capa[w]$ está definida.) El algoritmo
vale tanto para grafos no dirigidos como dirigidos (en ese caso $\Adj[t]$ son los sucesores de $t$).

\subsection*{(b) Corrección}
Sea $\delta(w)=\dist(u,w)$. \textbf{Recurrencia:}
\[
N_u=1,\qquad N_w=\sum_{\substack{t:\ (t,w)\in E\\ \delta(t)=\delta(w)-1}} N_t\quad\text{para } w\neq u \text{ alcanzable},
\qquad N_w=0 \text{ si } w \text{ no es alcanzable.}
\]
\begin{lema}
La recurrencia es correcta.
\end{lema}
\begin{proof}
$N_u=1$: el único camino de largo $0$ de $u$ a $u$ es el trivial. Sea $w\ne u$ alcanzable, con
$\delta(w)=d\ge1$, y sea $\mathcal{P}_w$ el conjunto de caminos mínimos de $u$ a $w$. Partimos
$\mathcal{P}_w$ según el penúltimo vértice $t$ del camino. Para cada $t$ definimos la función
$\varphi_t$ que a un camino mínimo $Q$ de $u$ a $t$ le asocia $Q$ seguido de la arista $(t,w)$.
\begin{itemize}
\item Si $P\in\mathcal{P}_w$ tiene penúltimo vértice $t$, entonces $(t,w)\in E$ y el prefijo $Q$ de $P$
hasta $t$ tiene largo $d-1$. $Q$ es un camino mínimo a $t$ (si hubiera uno más corto, agregándole
$(t,w)$ obtendríamos un recorrido de largo $<d$ de $u$ a $w$, absurdo), así que $\delta(t)=d-1$ y
$P=\varphi_t(Q)$.
\item Recíprocamente, si $(t,w)\in E$, $\delta(t)=d-1$ y $Q$ es un camino mínimo a $t$, entonces
$\varphi_t(Q)$ es un camino (no repite vértices: los vértices de $Q$ están a distancia $\le d-1$ de $u$
y $w$ a distancia $d$) de largo $d=\delta(w)$, o sea un camino mínimo a $w$ con penúltimo vértice $t$.
\end{itemize}
Cada $\varphi_t$ es inyectiva y caminos con distinto penúltimo vértice son distintos, así que
$|\mathcal{P}_w|=\sum_t N_t$ sobre los $t$ indicados. Si $w$ no es alcanzable no hay caminos: $N_w=0$.
\end{proof}

\textbf{El algoritmo calcula la recurrencia.} Por (3.3), $L_i=\{w:\delta(w)=i\}$ y $capa[w]=\delta(w)$
para todo $w$ descubierto. Probamos por inducción en $i$ el invariante: \emph{al comenzar la iteración $i$
del \textbf{mientras}, $N[w]=N_w$ para todo $w\in L_0\cup\dots\cup L_i$.}
\begin{itemize}
\item $i=0$: $N[u]=1=N_u$.
\item Paso: sea $w\in L_{i+1}$. Su valor $N[w]$ vale 0 hasta que se descubre, lo cual ocurre durante la
iteración $i$ (no antes, pues entonces estaría en una capa $\le i$). Solo se modifica en la línea 12
cuando se examina una arista $(t,w)$ con $t\in L_j$ y $capa[w]=j+1$, es decir con $j=i$. Durante la
iteración $i$ se recorre cada $t\in L_i$ exactamente una vez, y cada arista $(t,w)$ aparece una sola
vez en $\Adj[t]$; para cada una de ellas, $w$ queda descubierto (en esa misma vuelta si no lo estaba) con
$capa[w]=i+1$, y se suma $N[t]=N_t$ (hipótesis inductiva). Así, al terminar la iteración $i$,
\[N[w]=\sum_{t\in L_i,\ (t,w)\in E}N_t=\sum_{(t,w)\in E,\ \delta(t)=\delta(w)-1}N_t=N_w,\]
y este valor ya no cambia en iteraciones posteriores (para $j>i$, $capa[w]=i+1\ne j+1$). Los vértices de
capas $\le i$ tampoco cambian: para ellos $capa[w]\le i<j+1$ en toda iteración $j\ge i$.
\end{itemize}
Al terminar, los vértices descubiertos son exactamente los alcanzables y tienen $N[w]=N_w$; los no
alcanzables mantienen $N[w]=0=N_w$. En particular se retorna $N[v]=N_v$.

\subsection*{(c) Tiempo de ejecución}
Es BFS (con listas de adyacencia y las capas como listas enlazadas) más un test y una suma $\Oh(1)$ por
cada elemento de cada lista de adyacencia recorrida. Cada vértice está en una sola capa, por lo que
cada $\Adj[t]$ se recorre una sola vez: el costo es $\Oh(n+\sum_t\deg(t))=\Oh(n+m)$, más la
inicialización $\Oh(n)$. Total $\Oh(n+m)$, asumiendo (como es usual en el curso) que las sumas de
enteros cuestan $\Oh(1)$.\footnote{Las cantidades $N_w$ pueden crecer exponencialmente en $n$; con
aritmética de precisión arbitraria el costo de cada suma dejaría de ser constante.}

%======================================================================
\section{Ejercicio 7}
\begin{enunciado}
(a) Proponga un problema de la realidad que se modele con un grafo $G$ y cuya solución sea equivalente
a determinar si $G$ es bipartito.
(b) Demuestre que si $G$ tiene un ciclo con cantidad impar de aristas, entonces no es bipartito.
(c) Escriba un algoritmo basado en DFS que determine si un grafo es bipartito.
\end{enunciado}
\examref{Examen julio 2025, Ej.~1 (determinar con DFS si se puede repartir $V$ en dos conjuntos sin aristas internas y si la repartición es única); Examen julio 2016, Problema B (paridad de distancias y aristas entre vértices de igual paridad).}

\begin{definicion}
$G=(V,E)$ es \emph{bipartito} si existe una partición $V=X\cup Y$, $X\cap Y=\emptyset$, tal que toda
arista tiene un extremo en $X$ y otro en $Y$. Equivalentemente, existe una coloración $c:V\to\{0,1\}$
con $c(x)\neq c(y)$ para toda $\{x,y\}\in E$.
\end{definicion}

\subsection*{(a) Un problema de la realidad}
Un club quiere repartir a sus $n$ socios en dos turnos (mañana y tarde) para usar el gimnasio. Hay
pares de socios que están peleados y no pueden coincidir en el mismo turno. Modelo: un vértice por
socio, y una arista $\{x,y\}$ por cada par peleado. Un reparto válido es exactamente una partición de
$V$ en dos conjuntos (los turnos) donde cada arista une socios de turnos distintos: existe reparto
si y solo si $G$ es bipartito, y una bipartición es un reparto.

\subsection*{(b) Ciclo impar $\Rightarrow$ no bipartito}
Sea $C=v_1,v_2,\dots,v_k,v_1$ un ciclo con $k$ impar, y supongamos por absurdo que $c$ es una
2-coloración válida. Probamos por inducción que $c(v_i)=c(v_1)$ si $i$ es impar y $c(v_i)\ne c(v_1)$ si
$i$ es par: vale para $i=1$; y como $\{v_i,v_{i+1}\}\in E$, $c(v_{i+1})\ne c(v_i)$, es decir, el color
alterna en cada paso, así que la propiedad pasa de $i$ a $i+1$. Como $k$ es impar, $c(v_k)=c(v_1)$;
pero $\{v_k,v_1\}\in E$ exige $c(v_k)\neq c(v_1)$. Absurdo; $G$ no es bipartito.

\subsection*{(c) Algoritmo basado en DFS}
Coloreamos durante el DFS: la raíz de cada árbol recibe el color 0 y cada hijo el color opuesto al de
su padre. Al examinar una arista hacia un vértice ya explorado se controla que los colores difieran.

\begin{algorithm}[H]
\caption{Test de bipartición con DFS}
\begin{algorithmic}[1]
\Function{EsBipartito}{$G$}
  \State $color[v]\gets\textbf{nil}$ para todo $v\in V$
  \ForAll{$s\in V$}
    \If{$color[s]=\textbf{nil}$}
      \State $color[s]\gets0$
      \If{\textbf{not} \Call{DFS-Color}{$s$}} \State \Return \textbf{false} \EndIf
    \EndIf
  \EndFor
  \State \Return \textbf{true} \Comment{$X=\{v:color[v]=0\}$, $Y=\{v:color[v]=1\}$}
\EndFunction
\Statex
\Function{DFS-Color}{$u$}
  \ForAll{$v\in\Adj[u]$}
    \If{$color[v]=\textbf{nil}$}
      \State $color[v]\gets1-color[u]$
      \If{\textbf{not} \Call{DFS-Color}{$v$}} \State \Return \textbf{false} \EndIf
    \ElsIf{$color[v]=color[u]$}
      \State \Return \textbf{false} \Comment{hay un ciclo impar}
    \EndIf
  \EndFor
  \State \Return \textbf{true}
\EndFunction
\end{algorithmic}
\end{algorithm}
($color[v]\neq\textbf{nil}$ cumple el papel de $Explorado[v]$; el algoritmo es el DFS del libro sobre
todas las componentes.)

\textbf{Corrección.} En todo momento, para cada vértice coloreado $v$, $color[v]\equiv prof(v)\pmod 2$,
donde $prof$ es la profundidad en el árbol DFS de su componente (la raíz tiene color 0 y profundidad 0,
y cada hijo tiene profundidad $+1$ y color opuesto).
\begin{itemize}
\item \emph{Si retorna \textbf{true}, $G$ es bipartito.} Cada vértice se colorea (el bucle de la línea
3 lanza un DFS en cada componente), y cada \textsc{DFS-Color}($u$) se ejecuta una vez y examina todas
las aristas $\{u,v\}$. En ese momento, o bien $v$ no tenía color y recibe $1-color[u]$, o bien tenía
color y la línea 15 verifica que sea distinto del de $u$. Como no se retornó \textbf{false}, toda arista
tiene extremos de colores distintos y los colores no se modifican después de asignados: $(X,Y)$ es una
bipartición.
\item \emph{Si retorna \textbf{false}, $G$ no es bipartito.} Se encontró $\{u,v\}\in E$ con
$color[u]=color[v]$. No es arista de árbol (las de árbol unen padre e hijo, de colores distintos), así
que por (3.7) uno es ancestro del otro, digamos $v$ ancestro de $u$. Como tienen igual color,
$prof(u)-prof(v)$ es par y positivo, por lo que el camino del árbol de $v$ a $u$ tiene una cantidad par
$\ge2$ de aristas; agregando $\{u,v\}$ se obtiene un ciclo con una cantidad impar de aristas, y por (b)
$G$ no es bipartito.
\end{itemize}

\textbf{Orden.} Es un DFS completo con trabajo $\Oh(1)$ adicional por arista examinada: $\Oh(n+m)$.

%======================================================================
\section{Ejercicio 8}
\begin{enunciado}
(K\&T 3.4) Un coleccionista tiene $n$ mariposas de dos especies. Para algunos pares establece un juicio
y etiqueta el par como \emph{iguales} o \emph{diferentes} ($m$ juicios en total). Quiere saber si los
juicios son consistentes, es decir, si se puede asignar a cada mariposa una especie, $A$ o $B$, de
modo que todo par \emph{iguales} quede en la misma especie y todo par \emph{diferentes} en especies
distintas; y en ese caso obtener una clasificación.
(a) Modele el problema con conceptos matemáticos.
(b) Si todos los juicios son \emph{diferentes}, ¿es posible que sean inconsistentes? Si no, explique por
qué; si sí, escriba un algoritmo que resuelva el problema.
(c) Para el caso general escriba un algoritmo que resuelva el problema.
Todos los algoritmos deben ser $\Oh(m+n)$; reescriba los algoritmos del libro que utilice.
(d) Demuestre la corrección de sus algoritmos.
(e) Demuestre que admiten una implementación $\Oh(m+n)$.
\end{enunciado}
\examref{Examen diciembre 2025, Ej.~2: mismo problema, resuelto con un ``BFS modificado'' que pone en la misma capa los vértices unidos por aristas \emph{iguales} y en la capa siguiente los unidos por \emph{diferentes}, coloreando por paridad de capa y verificando luego todas las aristas. Nuestra solución es equivalente (la paridad de la capa es nuestro color).}

\subsection*{(a) Modelo}
Grafo no dirigido $G=(V,E)$ con $V$ = las $n$ mariposas y una arista $\{x,y\}$ por cada juicio, con una
etiqueta $\ell(x,y)\in\{\text{iguales},\text{diferentes}\}$; $|E|=m$. Codificamos las etiquetas con
$\delta(x,y)=0$ si es \emph{iguales} y $\delta(x,y)=1$ si es \emph{diferentes}, y las especies con
$A=0$, $B=1$. Una \emph{clasificación} es una función $c:V\to\{0,1\}$; la arista $\{x,y\}$ es
\emph{válida} en $c$ si
\[ c(x)\oplus c(y)=\delta(x,y)\qquad(\oplus=\text{suma módulo 2}), \]
es decir, colores iguales para \emph{iguales} y distintos para \emph{diferentes}. Los juicios son
\emph{consistentes} si existe una clasificación en la que todas las aristas son válidas
(\emph{clasificación válida}). El problema: decidir si existe y, en caso afirmativo, exhibir una.

\subsection*{(b) Todos los juicios \emph{diferentes}}
\textbf{Sí, pueden ser inconsistentes.} Con tres mariposas $1,2,3$ y los juicios ``1 y 2 diferentes'',
``2 y 3 diferentes'', ``1 y 3 diferentes'': si $c(1)=A$ entonces $c(2)=B$, $c(3)=A$, pero 1 y 3 deberían
ser diferentes. En general, con todas las etiquetas \emph{diferentes}, una clasificación válida es
exactamente una bipartición de $G$: los juicios son consistentes si y solo si $G$ es bipartito (y el
ejemplo es un ciclo impar, ver Ejercicio~7(b)). Algoritmo: el test de bipartición del libro con BFS.

\begin{algorithm}[H]
\caption{Juicios todos \emph{diferentes}: test de bipartición con BFS}
\begin{algorithmic}[1]
\Function{Bipartito}{$G$}
  \State $c[v]\gets\textbf{nil}$ para todo $v\in V$
  \ForAll{$s\in V$ con $c[s]=\textbf{nil}$} \Comment{un BFS por componente conexa}
    \State $c[s]\gets 0$; \ $L_0\gets\{s\}$; \ $i\gets0$
    \While{$L_i\ne\emptyset$}
      \State $L_{i+1}\gets\emptyset$
      \ForAll{$t\in L_i$, \ $w\in\Adj[t]$ con $c[w]=\textbf{nil}$}
        \State $c[w]\gets 1-c[t]$; \ agregar $w$ a $L_{i+1}$ \Comment{$c[w]$ = paridad de la capa}
      \EndFor
      \State $i\gets i+1$
    \EndWhile
  \EndFor
  \ForAll{$\{x,y\}\in E$}
    \If{$c[x]=c[y]$} \State \Return ``inconsistente'' \EndIf
  \EndFor
  \State \Return $c$ \Comment{especie $A$ si $c=0$, $B$ si $c=1$}
\EndFunction
\end{algorithmic}
\end{algorithm}
Es exactamente el caso particular $\delta\equiv1$ del algoritmo de la parte (c) (la línea 8 coincide con
$c[w]\gets c[t]\oplus\delta(t,w)$), por lo que su corrección y su orden se deducen de (d) y (e).

\subsection*{(c) Caso general}
Idea: propagar la especie por un BFS (uno por componente): la raíz recibe $A$ y cada vértice
descubierto por la arista $(t,w)$ recibe $c[w]=c[t]\oplus\delta(t,w)$, de modo que todas las aristas de
árbol son válidas. Al final se verifica que todas las aristas sean válidas.

\begin{algorithm}[H]
\caption{Consistencia de los juicios (caso general)}
\begin{algorithmic}[1]
\Function{Clasificar}{$G,\delta$}
  \State $c[v]\gets\textbf{nil}$ para todo $v\in V$
  \ForAll{$s\in V$ con $c[s]=\textbf{nil}$} \Comment{un BFS por componente conexa}
    \State $c[s]\gets0$; \ $L_0\gets\{s\}$; \ $i\gets0$
    \While{$L_i\neq\emptyset$}
      \State $L_{i+1}\gets\emptyset$
      \ForAll{$t\in L_i$}
        \ForAll{$w\in\Adj[t]$}
          \If{$c[w]=\textbf{nil}$}
            \State $c[w]\gets c[t]\oplus\delta(t,w)$; \ agregar $w$ a $L_{i+1}$
          \EndIf
        \EndFor
      \EndFor
      \State $i\gets i+1$
    \EndWhile
  \EndFor
  \ForAll{$\{x,y\}\in E$}
    \If{$c[x]\oplus c[y]\neq\delta(x,y)$} \State \Return ``inconsistente'' \EndIf
  \EndFor
  \State \Return $c$
\EndFunction
\end{algorithmic}
\end{algorithm}
La verificación se hace al final (y no al terminar cada componente) solo por simplicidad; como las
aristas no cruzan componentes, daría igual.

\subsection*{(d) Corrección}
Sea $F$ el conjunto de aristas de árbol, es decir las aristas $\{t,w\}$ por las cuales se descubrió $w$
en la línea 10. $F$ es un bosque con un árbol por componente conexa de $G$ (los árboles BFS).

\begin{lema}\label{l:arbol-valido}
Al terminar el bucle de la línea 3, todo vértice tiene $c[v]\in\{0,1\}$ y toda arista de $F$ es válida
en $c$.
\end{lema}
\begin{proof}
El bucle de la línea 3 lanza un BFS desde cada vértice aún sin color, y BFS colorea toda la componente
de $s$; así todos quedan coloreados. Un color, una vez asignado, no cambia (solo se asigna si era
\textbf{nil}). Si $\{t,w\}\in F$ con $w$ descubierto desde $t$, se asignó $c[w]=c[t]\oplus\delta(t,w)$,
o sea $c[t]\oplus c[w]=\delta(t,w)$: la arista es válida.
\end{proof}

\begin{lema}\label{l:camino}
Si $P=x_0,x_1,\dots,x_k$ es un camino y $g$ es una clasificación en la que todas las aristas de $P$
son válidas, entonces $g(x_k)=g(x_0)\oplus\bigoplus_{j=0}^{k-1}\delta(x_j,x_{j+1})$.
\end{lema}
\begin{proof}
Inducción en $k$: para $k=0$ es trivial, y $g(x_{k})=g(x_{k-1})\oplus\delta(x_{k-1},x_k)$ por validez de
la última arista.
\end{proof}

\begin{teorema}
\textsc{Clasificar} retorna una clasificación si y solo si los juicios son consistentes, y en ese caso
la clasificación retornada es válida.
\end{teorema}
\begin{proof}
\emph{Si retorna $c$:} por el Lema~\ref{l:arbol-valido} $c$ es una clasificación, y como pasó el control de
las líneas 12--13, todas las aristas son válidas. Luego $c$ es válida y los juicios son consistentes.

\emph{Si retorna ``inconsistente'':} hay una arista $\{x,y\}$ con $c[x]\oplus c[y]\neq\delta(x,y)$.
Como $x,y$ son adyacentes están en la misma componente, por lo tanto en el mismo árbol de $F$ con raíz
$r$; sean $P_x$, $P_y$ los caminos del árbol de $r$ a $x$ y a $y$, y $\Delta_x,\Delta_y$ la suma módulo
2 de las etiquetas de sus aristas. Por el Lema~\ref{l:arbol-valido} esas aristas son válidas en $c$, y
por el Lema~\ref{l:camino}: $c[x]=c[r]\oplus\Delta_x$ y $c[y]=c[r]\oplus\Delta_y$, de donde
$c[x]\oplus c[y]=\Delta_x\oplus\Delta_y$.

Supongamos por absurdo que existe una clasificación válida $g$. En $g$ todas las aristas son válidas,
en particular las de $P_x$ y $P_y$; por el Lema~\ref{l:camino}, $g(x)\oplus g(y)=(g(r)\oplus\Delta_x)
\oplus(g(r)\oplus\Delta_y)=\Delta_x\oplus\Delta_y=c[x]\oplus c[y]\neq\delta(x,y)$. Entonces $\{x,y\}$
no es válida en $g$: absurdo. Por lo tanto los juicios son inconsistentes.
\end{proof}
Para la parte (b) ($\delta\equiv1$), el teorema dice que el algoritmo retorna una 2-coloración
propia si y solo si $G$ es bipartito; además, en el caso inconsistente, $P_x$, $P_y$ y $\{x,y\}$
contienen un ciclo impar (el resultado del libro (3.15)), aunque esto no es necesario para la
corrección.

\subsection*{(e) Tiempo de ejecución}
Representamos $G$ con listas de adyacencia en las que cada entrada guarda el vecino y la etiqueta
$\delta$ de la arista; se construyen a partir de los $m$ juicios en $\Oh(n+m)$ (crear $n$ listas vacías y
agregar cada juicio a dos listas en $\Oh(1)$).
\begin{itemize}[nosep]
\item Inicialización: $\Oh(n)$.
\item Bucle de la línea 3: es el cálculo de todas las componentes conexas con BFS (ver (3.11)): cada vértice
recibe color, entra a una única capa y su lista se recorre una sola vez, con trabajo $\Oh(1)$ por
entrada. Total $\Oh(n+\sum_t\deg t)=\Oh(n+m)$.
\item Verificación final: recorriendo las listas de adyacencia se examina cada arista (dos veces) con
costo $\Oh(1)$: $\Oh(n+m)$.
\end{itemize}
Por la regla de la suma el total es $\Oh(n+m)$, tanto para (b) como para (c).

%======================================================================
\section{Ejercicio 9}
\begin{enunciado}
Consideramos el grafo dirigido $G$ de la figura.
(a) Ejecute BFS desde el nodo A.
(b) Ejecute DFS desde el nodo B.
(c) A partir de lo anterior, ¿qué puede decir de las componentes fuertemente conexas de A y B?
(d) Ejecute DFS desde el nodo A en $\Grev$.
(e) Encuentre todas las componentes fuertemente conexas de $G$. Justifique.
\end{enunciado}
\examref{Primer parcial 2013, Ej.~1(b) y Examen julio 2016, Ej.~3: la componente fuertemente conexa de $u$ se obtiene intersecando lo alcanzable desde $u$ en $G$ con lo alcanzable desde $u$ en el grafo transpuesto ($\Grev$).}

\begin{center}
\begin{tikzpicture}[scale=1.2]
  \node[nodo] (H) at (0.6,3.4) {H};
  \node[nodo] (A) at (2,2.6) {A};
  \node[nodo] (F) at (0.3,1.9) {F};
  \node[nodo] (E) at (1.8,1.1) {E};
  \node[nodo] (B) at (0.2,0.3) {B};
  \node[nodo] (C) at (3.3,1.6) {C};
  \node[nodo] (D) at (2.8,0) {D};
  \node[nodo] (G) at (4.1,0.3) {G};
  \draw[dir] (H)--(A); \draw[dir] (A)--(F); \draw[dir] (A)--(E); \draw[dir] (A)--(C);
  \draw[dir] (F)--(H); \draw[dir] (F)--(B); \draw[dir] (B)--(E); \draw[dir] (E)--(F);
  \draw[dir] (C)--(D); \draw[dir] (D)--(G); \draw[dir] (G)--(C);
\end{tikzpicture}

\small Aristas: $H\!\to\!A$, $A\!\to\!F$, $A\!\to\!E$, $A\!\to\!C$, $F\!\to\!H$, $F\!\to\!B$,
$B\!\to\!E$, $E\!\to\!F$, $C\!\to\!D$, $D\!\to\!G$, $G\!\to\!C$.
\end{center}
Usamos las listas de adyacencia (aristas salientes) en orden alfabético:
$A\!:\!C,E,F$;\ $B\!:\!E$;\ $C\!:\!D$;\ $D\!:\!G$;\ $E\!:\!F$;\ $F\!:\!B,H$;\ $G\!:\!C$;\ $H\!:\!A$.

\subsection*{(a) BFS desde A}
\begin{itemize}[nosep]
\item $L_0=\{A\}$.
\item A descubre C, E, F: \ $L_1=\{C,E,F\}$.
\item C descubre D; E nada nuevo (F ya descubierto); F descubre B y H: \ $L_2=\{D,B,H\}$.
\item D descubre G; B nada (E descubierto); H nada (A descubierto): \ $L_3=\{G\}$.
\item G nada (C descubierto): $L_4=\emptyset$.
\end{itemize}
Árbol BFS: $A\to C$, $A\to E$, $A\to F$, $C\to D$, $F\to B$, $F\to H$, $D\to G$. \textbf{Desde A se
alcanzan los 8 vértices.}
\begin{center}
\begin{tikzpicture}[level distance=10mm,
  level 1/.style={sibling distance=22mm},
  level 2/.style={sibling distance=12mm},
  every node/.style={nodo}, edge from parent/.style={draw,dir}]
\node {A}
  child {node {C} child {node {D} child {node {G}}}}
  child {node {E}}
  child {node {F} child {node {B}} child {node {H}}};
\end{tikzpicture}
\end{center}

\subsection*{(b) DFS desde B}
\textsc{DFS}(B) $\to$ E $\to$ F; en F: B explorado, sigue H $\to$ \textsc{DFS}(H) $\to$ A; en A: C
$\to$ \textsc{DFS}(C) $\to$ D $\to$ G; G: C explorado, retorna. De vuelta en A: E y F explorados.
Se retrocede hasta B sin nada nuevo.

Árbol DFS: $B\to E\to F\to H\to A\to C\to D\to G$ (un camino). \textbf{Desde B se alcanzan los 8
vértices.}

\subsection*{(c) Componentes fuertemente conexas de A y B}
Por (a), B es alcanzable desde A (hay camino $A\to F\to B$); por (b), A es alcanzable desde B (camino
$B\to E\to F\to H\to A$). Luego A y B son \emph{mutuamente alcanzables} y, por definición,
\textbf{pertenecen a la misma componente fuertemente conexa}. (Por (3.17) del libro las componentes
fuertemente conexas forman una partición de $V$, así que la componente de A es la de B.) Además, como
ambos alcanzan a todos los vértices, la componente de A está formada por los vértices que \emph{también}
alcanzan a A, lo cual se calcula en la parte siguiente.

\subsection*{(d) DFS desde A en $\Grev$}
$\Grev$ tiene las aristas invertidas. Listas (alfabéticas): $A\!:\!H$;\ $B\!:\!F$;\ $C\!:\!A,G$;\
$D\!:\!C$;\ $E\!:\!A,B$;\ $F\!:\!A,E$;\ $G\!:\!D$;\ $H\!:\!F$.

\textsc{DFS}(A) $\to$ H $\to$ F; en F: A explorado, sigue E $\to$ \textsc{DFS}(E): A explorado, sigue B
$\to$ \textsc{DFS}(B): F explorado. Se retrocede hasta A, sin nada nuevo.

Árbol: $A\to H\to F\to E\to B$. En $\Grev$ desde A se alcanzan $\{A,H,F,E,B\}$, que son exactamente los
vértices que \emph{alcanzan a A en $G$}.

\subsection*{(e) Todas las componentes fuertemente conexas}
La componente de A es el conjunto de vértices mutuamente alcanzables con A:
\[ \mathrm{CFC}(A)=\{v: A\leadsto v\}\cap\{v: v\leadsto A\}
=V\cap\{A,B,E,F,H\}=\{A,B,E,F,H\}. \]
Los vértices restantes son C, D, G. Hay un ciclo $C\to D\to G\to C$, así que los tres son mutuamente
alcanzables y están en la misma componente; esa componente no contiene a ningún otro vértice porque
las únicas aristas que salen de $\{C,D,G\}$ son $C\to D$, $D\to G$, $G\to C$, así que desde C no se
alcanza nada fuera de $\{C,D,G\}$. Como las componentes particionan $V$:
\[ \textbf{CFC de } G:\quad \{A,B,E,F,H\}\quad\text{y}\quad\{C,D,G\}. \]

%======================================================================
\section{Ejercicio 10}
\begin{enunciado}
Consideramos el grafo $G$ de la figura.
(a) ¿Qué nodos no tienen aristas entrantes?
(b) ¿Qué nodos no tienen aristas entrantes luego de eliminar los nodos de la parte anterior?
(c) Repita (a) y (b) en $\Grev$.
(d) Usando lo anterior encuentre al menos cuatro órdenes topológicos distintos para $G$.
(e) ¿Son todos los órdenes topológicos posibles?
(f) Describa un algoritmo para encontrar un orden topológico y demuestre que existe una
implementación $\Oh(m+n)$.
\end{enunciado}
\examref{Primer parcial 2022, Ej.~3(b) (análisis detallado en $\Oh(m+n)$ del algoritmo de orden topológico por eliminación de fuentes); Primer parcial 2009, Ej.~3 y Primer parcial 2015, Ej.~2 (orden topológico con DFS en postorden invertido).}

\begin{center}
\begin{tikzpicture}[scale=1.2]
  \node[nodo] (A) at (2,3) {A};
  \node[nodo] (F) at (0.5,2) {F};
  \node[nodo] (C) at (3.5,2) {C};
  \node[nodo] (E) at (2,1.4) {E};
  \node[nodo] (B) at (0.6,0.3) {B};
  \node[nodo] (D) at (3.4,0.3) {D};
  \draw[dir] (A)--(F); \draw[dir] (A)--(E); \draw[dir] (A)--(C);
  \draw[dir] (C)--(E); \draw[dir] (C)--(D); \draw[dir] (E)--(F);
  \draw[dir] (B)--(F); \draw[dir] (B)--(E);
\end{tikzpicture}

\small Aristas: $A\!\to\!F$, $A\!\to\!E$, $A\!\to\!C$, $C\!\to\!E$, $C\!\to\!D$, $E\!\to\!F$, $B\!\to\!F$, $B\!\to\!E$.
\end{center}

\subsection*{(a) y (b)}
Grados de entrada: $A\!:\!0$, $B\!:\!0$, $C\!:\!1$ (de A), $D\!:\!1$ (de C), $E\!:\!3$ (de A, B, C), $F\!:\!3$
(de A, B, E).
\begin{enumerate}[label=(\alph*)]
\item Sin aristas entrantes: \textbf{A y B}.
\item Eliminando A y B: C pierde su única entrante; E conserva $C\to E$; D conserva $C\to D$; F conserva
$E\to F$. Sin entrantes: \textbf{C}. (Si se sigue el proceso: al eliminar C quedan sin entrantes D y E;
luego F.)
\end{enumerate}

\subsection*{(c) Lo mismo en $\Grev$}
En $\Grev$ las entrantes de un nodo son las salientes en $G$.
\begin{itemize}[nosep]
\item Sin entrantes en $\Grev$ (sumideros de $G$): \textbf{D y F}.
\item Eliminando D y F: E solo tenía saliente hacia F, así que queda sin entrantes en $\Grev$; C conserva
$C\to E$; A conserva $A\to C$, $A\to E$; B conserva $B\to E$. Sin entrantes: \textbf{E}. (Siguiendo: luego
B y C; por último A.)
\end{itemize}

\subsection*{(d) Cuatro órdenes topológicos}
Eliminando fuentes en $G$ (partes (a), (b)) el orden empieza con A y B en cualquier orden, sigue C, luego
D y E en cualquier orden, y termina en F; eliminando fuentes en $\Grev$ y leyendo al revés (parte (c))
se termina con D y F en cualquier orden, antes E, antes B y C en cualquier orden, y al principio A.
Por ejemplo:
\[
\begin{array}{ll}
1.\ A,B,C,D,E,F & 2.\ B,A,C,D,E,F\\
3.\ A,B,C,E,D,F & 4.\ B,A,C,E,D,F\\
5.\ A,C,B,E,F,D & 6.\ A,C,B,E,D,F
\end{array}
\]
Cada uno respeta las 8 aristas (todas van de izquierda a derecha).

\subsection*{(e) ¿Son todos?}
\textbf{No.} Las aristas imponen $A<C<E<F$ (por $A\to C$, $C\to E$, $E\to F$), $B<E$ (y entonces
$B<F$), y $C<D$ (y entonces $A<D$); $A\to E$, $A\to F$ y $B\to F$ son consecuencia de las anteriores. Un
orden topológico es entonces la cadena $A,C,E,F$ con B insertado en uno de los 3 huecos anteriores a E
(antes de A, entre A y C, o entre C y E) y D insertado en uno de los 3 huecos posteriores a C (entre C y
E, entre E y F, o después de F). Son $3\cdot3=9$ elecciones de huecos, y cuando ambos caen en el hueco
entre C y E hay 2 órdenes (B antes o después de D). En total \textbf{10 órdenes topológicos}:
\[
\begin{array}{lllll}
BACDEF & BACEDF & BACEFD & ABCDEF & ABCEDF\\
ABCEFD & ACBDEF & ACDBEF & ACBEDF & ACBEFD
\end{array}
\]
(por ejemplo $A,C,D,B,E,F$ no aparece en la parte (d)).

\subsection*{(f) Algoritmo y orden $\Oh(m+n)$}
Algoritmo del libro: por (3.19) todo DAG tiene un nodo sin aristas entrantes; se lo pone primero, se lo
elimina (el resto sigue siendo un DAG) y se repite. Para hacerlo eficiente mantenemos $cnt[w]$ = cantidad
de aristas entrantes a $w$ desde nodos no eliminados, y el conjunto $S$ de nodos no eliminados con
$cnt=0$.

\begin{algorithm}[H]
\caption{Orden topológico por eliminación de fuentes}
\begin{algorithmic}[1]
\Function{OrdenTopologico}{$G$} \Comment{$G$ dirigido acíclico}
  \State $cnt[v]\gets0$ para todo $v$
  \ForAll{$u\in V$, $w\in\Adj[u]$} \State $cnt[w]\gets cnt[w]+1$ \EndFor
  \State $S\gets\{v: cnt[v]=0\}$ \Comment{pila o cola}
  \State $orden\gets\langle\rangle$
  \While{$S\neq\emptyset$}
    \State sacar un $v$ de $S$; \ agregar $v$ al final de $orden$
    \ForAll{$w\in\Adj[v]$} \Comment{eliminar $v$}
      \State $cnt[w]\gets cnt[w]-1$
      \If{$cnt[w]=0$} agregar $w$ a $S$ \EndIf
    \EndFor
  \EndWhile
  \State \Return $orden$
\EndFunction
\end{algorithmic}
\end{algorithm}

\textbf{Corrección} (resumen; la demostración completa, que además trata grafos con ciclos, está en el
Ejercicio~11): se mantiene el invariante ``$cnt[w]$ es la cantidad de predecesores de $w$ no
eliminados''. Cuando $v$ se agrega a $orden$, $cnt[v]=0$, o sea todos sus predecesores ya están en
$orden$: toda arista $(u,v)$ va de un nodo anterior a uno posterior. Y si $G$ es un DAG, mientras queden
nodos no eliminados, el subgrafo que inducen es un DAG y por (3.19) tiene una fuente, que está en $S$;
así se eliminan los $n$ nodos.

\textbf{Orden.} Con listas de adyacencia:
\begin{itemize}[nosep]
\item Inicializar $cnt$ (líneas 2--4) recorre todas las listas: $\Oh(n+m)$; construir $S$: $\Oh(n)$.
\item Con $S$ implementado como pila o cola (lista enlazada) insertar y sacar cuesta $\Oh(1)$. Cada nodo
entra a $S$ a lo sumo una vez ($cnt[w]$ solo decrece, así que llega a 0 una sola vez; o era 0 desde el
inicio y nunca se decrementa), por lo que el \textbf{mientras} itera a lo sumo $n$ veces.
\item En la iteración que procesa $v$ el \textbf{para} interno recorre $\Adj[v]$ con trabajo $\Oh(1)$ por
arista: en total $\sum_v|\Adj[v]|=m$.
\end{itemize}
Por la regla de la suma: $\Oh(n+m)$.

%======================================================================
\section{Ejercicio 11}
\begin{enunciado}
(K\&T 3.3) Dado un grafo dirigido arbitrario $G$ (no necesariamente un DAG) se quiere obtener un orden
topológico de $G$ o un ciclo de $G$.
(a) Presente un algoritmo que resuelve el problema en $\Oh(m+n)$ (reescriba los algoritmos del libro).
\emph{Sugerencia:} use las ideas de la demostración de (3.19).
(b) Demuestre la corrección.
(c) Demuestre que admite una implementación $\Oh(m+n)$.
\end{enunciado}
\examref{Examen diciembre 2024, Ej.~1: se adapta el algoritmo de orden topológico para que, si en algún momento no hay vértices con grado de entrada 0, informe que no hay solución.}

\subsection*{(a) Algoritmo}
Primera fase: el algoritmo de eliminación de fuentes del Ejercicio~10(f). Si elimina todos los nodos,
devuelve el orden. Si no, los nodos no eliminados forman un subgrafo $H$ en el que todo nodo tiene una
arista entrante desde $H$. Segunda fase (idea de la prueba de (3.19)): desde cualquier nodo de $H$
caminamos \emph{hacia atrás} por aristas entrantes dentro de $H$; como $H$ es finito, en algún momento
repetimos un nodo, y el tramo entre las dos apariciones es un ciclo.

\begin{algorithm}[H]
\caption{Orden topológico o ciclo}
\begin{algorithmic}[1]
\Function{OrdenOCiclo}{$G$}
  \State Construir las listas de entrantes $In[v]=\{u:(u,v)\in E\}$
  \State $cnt[v]\gets|In[v]|$; \ $elim[v]\gets\textbf{false}$ para todo $v$
  \State $S\gets\{v:cnt[v]=0\}$; \ $orden\gets\langle\rangle$
  \While{$S\ne\emptyset$} \Comment{Fase 1}
    \State sacar $v$ de $S$; \ agregar $v$ a $orden$; \ $elim[v]\gets\textbf{true}$
    \ForAll{$w\in\Adj[v]$}
      \State $cnt[w]\gets cnt[w]-1$
      \If{$cnt[w]=0$} agregar $w$ a $S$ \EndIf
    \EndFor
  \EndWhile
  \If{$|orden|=n$} \State \Return ``orden topológico'', $orden$ \EndIf
  \Statex \Comment{Fase 2: caminata hacia atrás dentro de $H=\{v:\neg elim[v]\}$}
  \State $pos[v]\gets-1$ para todo $v$
  \State elegir $x$ con $\neg elim[x]$; \ $k\gets0$; \ $w_0\gets x$; \ $pos[x]\gets0$
  \Loop
    \State buscar en $In[w_k]$ un $u$ con $\neg elim[u]$
    \If{$pos[u]\ne-1$} \Comment{$u$ ya fue visitado: $u=w_j$ con $j=pos[u]$}
      \State $j\gets pos[u]$
      \State \Return ``ciclo'', $\langle w_k, w_{k-1},\dots,w_{j}\rangle$ \Comment{cerrado por la arista $(w_j,w_k)$}
    \EndIf
    \State $k\gets k+1$; \ $w_k\gets u$; \ $pos[u]\gets k$
  \EndLoop
\EndFunction
\end{algorithmic}
\end{algorithm}

\subsection*{(b) Corrección}
\textbf{Invariante de la fase 1:} $cnt[w]=|\{u\in In[w]:\neg elim[u]\}|$ para todo $w$, y $S$ contiene
exactamente a los nodos no eliminados con $cnt=0$ que aún no fueron sacados. Vale al inicio; al eliminar
$v$ se decrementa $cnt[w]$ una vez por cada arista $(v,w)$, que es exactamente lo que $v$ aportaba; y un
nodo se agrega a $S$ justo cuando su $cnt$ pasa a valer 0 (como $cnt$ solo decrece, eso ocurre a lo sumo
una vez, y los nodos con $cnt=0$ inicial no se decrementan nunca), así que cada nodo entra a $S$ a lo
sumo una vez.

\textbf{Si retorna un orden, es topológico.} Si $v$ se agrega a $orden$ cuando $cnt[v]=0$, por el
invariante todos sus predecesores $u$ ya estaban eliminados, es decir, ya estaban antes en $orden$. Luego
para toda arista $(u,v)$, $u$ aparece antes que $v$; y como $|orden|=n$ y ningún nodo se repite, es un
orden topológico de $G$.

\textbf{Si se llega a la fase 2, la caminata está bien definida.} Al terminar la fase 1, $S=\emptyset$
y $|orden|<n$, así que $H\ne\emptyset$ y todo $v\in H$ tiene $cnt[v]\ge1$ (si fuera 0 habría entrado a
$S$ y habría sido eliminado). Por el invariante, cada $v\in H$ tiene un predecesor $u\in H$: la búsqueda
de la línea 15 siempre encuentra uno, y la caminata permanece en $H$.

\textbf{La fase 2 termina y retorna un ciclo.} Mientras no se repite un nodo, los $w_0,\dots,w_k$ son
distintos, así que $k<|H|\le n$; por lo tanto en a lo sumo $|H|$ pasos el $u$ encontrado ya fue visitado
($pos[u]\neq-1$), digamos $u=w_j$ con $j\le k$. Por construcción $(w_{i+1},w_i)\in E$ para $0\le i<k$
(cada $w_{i+1}$ se eligió en $In[w_i]$) y además $(w_j,w_k)=(u,w_k)\in E$. Entonces la secuencia
\[ w_k\to w_{k-1}\to\dots\to w_{j+1}\to w_j\to w_k \]
usa solo aristas de $E$ y sus nodos $w_j,\dots,w_k$ son distintos: es un ciclo dirigido de $G$ (de
largo $k-j+1\ge1$; si $G$ no tiene lazos, $u\ne w_k$ y el largo es $\ge2$).

\textbf{Conclusión.} El algoritmo siempre termina con uno de los dos resultados y el resultado es
correcto. (En particular, retorna un orden topológico si y solo si $G$ es un DAG: un grafo con ciclos no
tiene orden topológico (3.18), y un DAG no tiene ciclos para devolver.)

\subsection*{(c) Tiempo de ejecución}
Con listas de adyacencia de salientes y de entrantes:
\begin{itemize}[nosep]
\item Construir $In$ recorriendo todas las listas $\Adj$: $\Oh(n+m)$. Inicializar $cnt$, $elim$, $S$, $pos$: $\Oh(n)$.
\item Fase 1: como en el Ejercicio~10(f), cada nodo sale de $S$ a lo sumo una vez y su lista de salientes
se recorre una vez con trabajo $\Oh(1)$ por arista: $\Oh(n+m)$.
\item Fase 2: cada nodo aparece como $w_k$ a lo sumo una vez (los $w_i$ son distintos hasta la
repetición, en cuyo momento se termina), y para $w_k$ la búsqueda de la línea 15 recorre a lo sumo
$|In[w_k]|$ entradas con costo $\Oh(1)$. Total $\Oh(\sum_v(1+|In[v]|))=\Oh(n+m)$. Para devolver el
ciclo guardamos los $w_i$ en un arreglo indexado por $i$; listar $w_k,\dots,w_j$ cuesta $\Oh(n)$.
\end{itemize}
Por la regla de la suma, el total es $\Oh(n+m)$.

%======================================================================
\section{Ejercicio 12}
\begin{enunciado}
(K\&T 3.12) Se tienen $n$ personas ya fallecidas $P_1,\dots,P_n$ y $m$ datos sobre pares $(P_i,P_j)$ de
alguno de los tipos: (1) $P_i$ murió antes de que $P_j$ naciera; (2) las vidas de $P_i$ y $P_j$ se
solaparon en el tiempo. Se quiere un algoritmo que devuelva una cronología de nacimientos y muertes de
las $n$ personas consistente con todos los datos (respetando que cada persona nace antes de morir), o
un mensaje indicando que los datos son inconsistentes. Por ejemplo, si $n=3$, $P_1$ murió antes de que
$P_3$ naciera y las vidas de $P_2$ y $P_1$ se solaparon, una salida posible es: nace $P_2$, nace $P_1$,
muere $P_1$, nace $P_3$, muere $P_2$, muere $P_3$.
(a) Presente un algoritmo $\Oh(m+n)$; describa la representación de la entrada y formalice un modelo.
(b) Demuestre la corrección.
(c) Demuestre que admite una implementación $\Oh(m+n)$.
\end{enunciado}
\examref{Examen diciembre 2024, Ej.~1: problema prácticamente idéntico (tareas con restricciones ``$j$ empieza después de que termina $i$'' ($T_1$) y ``los lapsos de $i$ y $j$ tienen un instante común'' ($T_2$)); la solución oficial construye el mismo grafo de $2n$ vértices y busca un orden topológico. Seguimos esa solución.}

\subsection*{(a) Entrada, modelo y algoritmo}
\textbf{Entrada.} El entero $n$ y una lista de $m$ datos, cada uno una terna $(i,j,\tau)$ con
$i,j\in\{1,\dots,n\}$ y $\tau\in\{\textsc{Antes},\textsc{Solapa}\}$: $(i,j,\textsc{Antes})$ significa
``$P_i$ murió antes de que $P_j$ naciera'' y $(i,j,\textsc{Solapa})$ ``las vidas de $P_i$ y $P_j$ se
solaparon''.

\textbf{Modelo.} Cada persona tiene dos eventos: su nacimiento $b_i$ y su muerte $d_i$. Una
\emph{cronología} es una secuencia de los $2n$ eventos (suponemos, como en el ejemplo, que no hay dos
eventos simultáneos; la cronología es el orden en que ocurren). Una cronología es \emph{consistente} si:
\begin{enumerate}[nosep,label=(R\arabic*)]
\item $b_i$ precede a $d_i$ para todo $i$;
\item para cada dato $(i,j,\textsc{Antes})$, $d_i$ precede a $b_j$;
\item para cada dato $(i,j,\textsc{Solapa})$, los intervalos de vida $[b_i,d_i]$ y $[b_j,d_j]$ tienen un
instante en común. Dado (R1), esto equivale a que \emph{$b_i$ preceda a $d_j$ y $b_j$ preceda a $d_i$}:
si comparten un instante $x$, entonces $b_i\le x\le d_j$ y $b_j\le x\le d_i$ (y las desigualdades
entre eventos distintos son estrictas); recíprocamente, si $b_i<d_j$ y $b_j<d_i$, el instante
$\max(b_i,b_j)$ es menor que $d_i$ y que $d_j$ (usando además (R1)), así que pertenece a ambas vidas.
\end{enumerate}
Todas las condiciones son precedencias entre pares de eventos. Construimos el grafo dirigido $G=(V,E)$ con
$V=\{b_1,d_1,\dots,b_n,d_n\}$ y las aristas:
\begin{itemize}[nosep]
\item $(b_i,d_i)$ para cada $i$;
\item $(d_i,b_j)$ para cada dato $(i,j,\textsc{Antes})$;
\item $(b_i,d_j)$ y $(b_j,d_i)$ para cada dato $(i,j,\textsc{Solapa})$.
\end{itemize}
Por construcción: \emph{una secuencia de los $2n$ eventos es una cronología consistente si y solo si es un
orden topológico de $G$} (cada arista representa exactamente una de las precedencias exigidas, y cada
precedencia exigida tiene su arista).

\textbf{Algoritmo.}
\begin{algorithm}[H]
\caption{Cronología consistente}
\begin{algorithmic}[1]
\Function{Cronologia}{$n$, $D$}
  \State crear $G$ con vértices $b_1,d_1,\dots,b_n,d_n$ y listas de adyacencia vacías
  \For{$i\gets1$ \textbf{hasta} $n$} agregar la arista $(b_i,d_i)$ \EndFor
  \ForAll{$(i,j,\tau)\in D$}
    \If{$\tau=\textsc{Antes}$} agregar $(d_i,b_j)$
    \Else\ agregar $(b_i,d_j)$ y $(b_j,d_i)$
    \EndIf
  \EndFor
  \State $(res, sec)\gets$ \Call{OrdenOCiclo}{$G$} \Comment{Ejercicio 11}
  \If{$res=$ ``orden topológico''}
    \State \Return $sec$, leyendo $b_i$ como ``nace $P_i$'' y $d_i$ como ``muere $P_i$''
  \Else
    \State \Return ``los datos son inconsistentes''
  \EndIf
\EndFunction
\end{algorithmic}
\end{algorithm}
En el ejemplo del enunciado las aristas son $(b_1,d_1)$, $(b_2,d_2)$, $(b_3,d_3)$, $(d_1,b_3)$, $(b_1,d_2)$, $(b_2,d_1)$,
y $b_2,b_1,d_1,b_3,d_2,d_3$ es un orden topológico: es la salida mostrada.

\subsection*{(b) Corrección}
\begin{itemize}
\item \textbf{Si el algoritmo devuelve una secuencia, es una cronología consistente.} Por el
Ejercicio~11(b), la secuencia es un orden topológico de $G$; por la observación del modelo, cumple (R1),
(R2) y (R3). Asignando a cada evento como instante su posición en la secuencia se obtienen $n$ personas
(cada una nace antes de morir) que satisfacen todos los datos.
\item \textbf{Si los datos son consistentes, el algoritmo devuelve una secuencia.} Si existió un conjunto
de $n$ personas que satisface todos los datos, ordenando sus $2n$ eventos por instante de ocurrencia se
obtiene una secuencia que cumple (R1)--(R3), es decir, un orden topológico de $G$. Por (3.18), $G$ es
entonces un DAG, así que no tiene ciclos y \textsc{OrdenOCiclo} (que por el Ejercicio~11 siempre
devuelve un orden topológico o un ciclo correcto) devuelve un orden topológico.
\end{itemize}
Equivalentemente: el algoritmo informa ``inconsistente'' exactamente cuando $G$ tiene un ciclo, y el ciclo
hallado es incluso una justificación de la inconsistencia (una cadena circular de precedencias).

\subsection*{(c) Tiempo de ejecución}
$G$ tiene $|V|=2n$ vértices y $|E|=n+m_1+2m_2\le n+2m$ aristas, donde $m_1$ y $m_2$ son las
cantidades de datos de cada tipo ($m_1+m_2=m$). Identificando $b_i$ con el entero $i$ y $d_i$ con $n+i$,
la construcción de las listas de adyacencia cuesta $\Oh(n)$ para crearlas y $\Oh(1)$ por arista
agregada: $\Oh(n+m)$. Por el Ejercicio~11(c), \textsc{OrdenOCiclo} cuesta $\Oh(|V|+|E|)=\Oh(2n+n+2m)=
\Oh(n+m)$, y traducir la secuencia a texto es $\Oh(n)$. Por la regla de la suma, el total es $\Oh(n+m)$.

\end{document}
